\documentclass[12pt]{article}
\usepackage{amsmath,amssymb}
\usepackage[margin=1in]{geometry}
\usepackage{hyperref}
\usepackage{xcolor}
\usepackage{booktabs}

\newcommand{\referee}[1]{\medskip\noindent\textcolor{gray}{\textit{Referee: ``#1''}}\medskip}
\newcommand{\response}[1]{\noindent\textbf{Response:} #1\medskip}

\begin{document}

\noindent
Dear Editor,

\medskip
We thank the referee for a careful and technically precise reading of our manuscript.
The report identifies several places where the exposition was insufficiently explicit,
and we have addressed each point. Following further independent mathematical review,
we have also corrected two additional technical errors in the paper: a mismatched
comparison baseline in Fig.~5 and an overstated endpoint divergence claim for the
single-mode metric formula. We respond to all issues below.

\bigskip\hrule\bigskip

%------- PART I -------
\section*{Part I: Referee Report Issues (1--6)}

%--- Issue 1 ---
\subsection*{Issue 1: Reality $\mathcal{R}$ is under-specified}
\referee{``The equivalence relation is not rigorously defined\ldots $\mathcal{R}$ risks
collapsing to a ray in Hilbert space.''}
\response{We have revised Definition~3 to state the equivalence relation explicitly:
$|\psi\rangle_F \sim |\psi\rangle_{F'}$ iff there exists $U_{F\to F'}$ with
$|\psi\rangle_{F'} = U_{F\to F'}|\psi\rangle_F$. We have added a sentence
clarifying that $\mathcal{R}$ carries the \emph{pointer structure} as additional
data---the specification of which observables behave classically. Two states related
by a unitary that does not preserve the pointer observable belong to different
equivalence classes, so $\mathcal{R}$ is a point on $M\times\Sigma$, not merely a
ray in $\mathcal{H}$. See revised Definition~3.}

%--- Issue 2 ---
\subsection*{Issue 2: $T_{AB}$ double-counts the K\"{a}hler structure}
\referee{``Your $T_{AB}$ is not new unless $\Omega_{AB}$ is not the standard symplectic
form. You don't clearly distinguish Berry curvature from new dynamical curvature.''}
\response{We have added a clarifying paragraph in Sec.~II.D. $\Omega_{AB}$ is
\emph{not} the Berry curvature: Berry curvature is holonomy of the state vector under
adiabatic parameter variation, living on the external parameter manifold. $\Omega_{AB}$
is curvature of the frame bundle over $\Sigma$, vanishing for reversible frame changes
and nonzero when a closed loop in $\Sigma$ fails to restore the pointer observable.
The symmetric $G_{AB}$ reproduces the Fubini--Study (K\"{a}hler) metric;
$\Omega_{AB}$ encodes dynamical irreversibility, absent from the standard K\"{a}hler
form. See the added paragraph following Definition~4.}

%--- Issue 3 ---
\subsection*{Issue 3: Axiom (A1) is equivalent to Gleason's non-contextuality}
\referee{``You have not eliminated the core assumption---you've rephrased it\ldots
[A1] is equivalent to non-contextuality across refinements, which is exactly what
Gleason assumes.''}
\response{We dispute the equivalence and have added an explicit comparison in
Sec.~IV.C. Gleason's non-contextuality (NC) is a \emph{global} condition: a
projector must yield the same probability across \emph{all} POVMs simultaneously.
Axiom~(A1) is strictly weaker: it constrains only isometric ancilla embeddings of
the form in Proposition~1. The logical relationship is NC $\Rightarrow$ (A1) but
(A1) $\not\Rightarrow$ NC. We additionally note that Gleason's theorem requires
$\dim\mathcal{H}\geq 3$, whereas our derivation holds for $\dim\geq 2$ because the
splitting operates on the ancilla-extended space. See the added paragraph in
Sec.~IV.C.}

%--- Issue 4 ---
\subsection*{Issue 4: Axiom (A3) is a strong, underived assumption}
\referee{``[A3] excludes contextual hidden-variable models\ldots another strong
assumption, not derived.''}
\response{The referee is correct, and we have revised Axiom~(A3) to label it
explicitly as a \emph{frame-locality postulate}---structurally analogous to local
causality in SR, adopted rather than derived. We explain why it is physically
motivated: frame-history dependence would undermine the coherence-frame programme
from the outset. See revised Theorem~2, Axiom~(A3).}

%--- Issue 5 ---
\subsection*{Issue 5: New physics claims contradict using standard QM structure}
\referee{``Your framework uses Hilbert space, unitary transformations, Born rule.
That is exactly standard quantum mechanics. Therefore you should get identical
predictions\ldots You claim both simultaneously, which is inconsistent.''}
\response{The source of the predictions is the \emph{geometric structure of $\Sigma$},
specifically the Fubini--Study metric $G_{\lambda\lambda}$, which standard QM does not
model. Standard QM implicitly treats all traversals of $\lambda$ as metrically
equivalent; Coherence Relativity weights physically traversed paths by
$G_{\lambda\lambda}$. No Hamiltonian is modified. We acknowledge that the geometric
structure of $\Sigma$ is additional structure beyond standard QM's formalism and have
revised Sec.~VI.D to describe the framework as a ``minimal geometric extension''
rather than a pure reinterpretation. See the added paragraph at the start of
Sec.~VI.D.}

%--- Issue 6 ---
\subsection*{Issues 6 \& 7: Hysteresis and $\tau_\mathrm{frame}$}
\referee{``In standard QM, evolution is unitary $\to$ path-independent probabilities.
No geometric hysteresis appears unless there is non-unitary dynamics. You haven't shown
where unitarity breaks.''}
\response{Unitarity does not break. We have added a sentence after Eq.~(13) making
this explicit: the full system--environment state evolves unitarily; the visibility
deficit is a property of the \emph{reduced} dynamics on $\Sigma$, the action cost
accumulated along the erasure--revival cycle under the effective stochastic drift.
Standard QM's path carries no metric weight; Coherence Relativity's does. Regarding
$\tau_\mathrm{frame}$: the ``Remark on universality'' box (p.~8) already states
explicitly that the timescales are not universal but depend on coupling strength
$v_\mathrm{max}$. The testable content is the ratio of transformation duration to
coupling strength, not the absolute value. No revision required for this sub-point.}

\bigskip\hrule\bigskip

%------- PART II -------
\section*{Part II: Additional Corrections (Independent Mathematical Review)}

Following independent review, we have identified and corrected two further technical
errors.

%--- Issue 7 (NEW) ---
\subsection*{Correction A (Critical): Fig.~5 comparison baseline}
\textbf{Error identified:} The paper defines $\lambda \equiv 1 - |\langle
W_L|W_R\rangle|$ but Fig.~5 plots $\mathcal{V}_\mathrm{std} = \sqrt{1-\lambda^2}$
as the standard QM reference curve. This formula is correct only for the orthogonal
distinguishability parameterization $D = \sqrt{1-|\langle W_L|W_R\rangle|^2}$.
Under the paper's own $\lambda$ definition, the correct standard QM prediction is
$\mathcal{V}_\mathrm{std} = 1 - \lambda$ (linear). The metric formula
$G_{\lambda\lambda} = 1/[2\lambda(2-\lambda)]$ is correctly derived; only the
comparison baseline in Fig.~5 was wrong. This error was verified numerically:
the spurious gap between $\sqrt{1-\lambda^2}$ and $1-\lambda$ reaches $0.37$ at
$\lambda = 0.5$, substantially larger than the claimed $\Delta \approx 0.24$.

\textbf{Correction:} We have redefined $\lambda$ as the standard
Englert--Greenberger--Yasin which-path distinguishability
$\lambda \equiv \sqrt{1-|\langle W_L|W_R\rangle|^2}$, which satisfies
$\mathcal{V}^2 + \lambda^2 = 1$ for pure states and makes
$\mathcal{V}_\mathrm{std} = \sqrt{1-\lambda^2}$ the correct baseline.
The metric $G_{\lambda\lambda}$ has been rederived under this parameterization,
and Fig.~5 and the $\Delta$ claim have been updated accordingly.

%--- Issue 8 (NEW) ---
\subsection*{Correction B: Metric endpoint divergence claim}
\textbf{Error identified:} The analytic formula $G_{\lambda\lambda} =
1/[2\lambda(2-\lambda)]$ approaches $1/2$ as $\lambda \to 1$ (finite, not divergent).
The claim ``the metric diverges at both endpoints'' was true only for the multi-mode
$N = 10$ numerical simulation in Fig.~1, not for the single-mode analytic formula.

\textbf{Correction:} Sec.~II.C and the Fig.~1 caption have been revised to state
that divergence at the classical endpoint $\lambda \to 1$ is a feature of the
multi-mode numerical model, not of the single-mode formula. The $\tau_\mathrm{frame}$
integral in Sec.~VI.D.1 uses the single-mode formula and correctly excludes the
coordinate-singular endpoints via the restricted range $\lambda \in [0.01, 0.99]$.

%--- Issues 9, 10 (NEW) ---
\subsection*{Corrections C \& D: Prediction provenance and $\Sigma$ dimension}
\textbf{C (Sec.~VI.D):} We have added a paragraph clarifying which predictions
follow from $G_{\lambda\lambda}$ alone (predictions~1 and~5: non-linear visibility
and transformation timescale) vs.\ which additionally invoke the deferred stochastic
dynamics from Paper~2 (predictions~2, 3, 4, 6). The quantitative falsifiability of
predictions~1 and~5 is independent of Paper~2.

\textbf{D (Sec.~III.C):} We have added one sentence after Proposition~1 clarifying
that frame splittings embed $\Sigma$ into a larger $\Sigma' = U(dM)/T^{dM}$, but
$T_{AB}$ need not be defined on $\Sigma'$ for the Born rule derivation; the splitting
is used solely to establish the probability assignment and the isometry returns to the
original $\Sigma$.

\bigskip\hrule\bigskip

\noindent\textbf{Summary of all changes:}

\medskip
\begin{tabular}{clp{7.5cm}}
\toprule
\# & Section & Change \\
\midrule
1 & Sec.~II.A Def.~3 & Added explicit equivalence relation; pointer structure clarified \\
2 & Sec.~II.D & Added paragraph: $\Omega_{AB}$ vs.\ Berry curvature \\
3 & Sec.~IV.C & Added paragraph: (A1) vs.\ Gleason NC; $\dim\geq 2$ note \\
4 & Sec.~IV Thm.~2 & Axiom~(A3) relabeled as frame-locality postulate \\
5 & Sec.~VI.D & Added paragraph: predictions from geometry, not modified $H$ \\
6 & Sec.~VI.D.1 & Added paragraph: hysteresis without unitarity violation \\
7 & Sec.~II.C, Fig.~5 & \textbf{[New]} Corrected $\lambda$ definition and comparison baseline \\
8 & Sec.~II.C, Fig.~1 & \textbf{[New]} Corrected endpoint divergence claim (single-mode formula) \\
9 & Sec.~VI.D & \textbf{[New]} Clarified which predictions require Paper~2 dynamics \\
10 & Sec.~III.C & \textbf{[New]} Added $\Sigma$-dimension note after Prop.~1 \\
\bottomrule
\end{tabular}

\medskip
\noindent We believe the revised manuscript fully addresses all raised concerns.
The core framework, Born rule derivation, and qualitative predictions are unchanged;
the revisions clarify foundations, correct two technical errors, and sharpen the
physical claims.

\bigskip
\noindent Sincerely,\\
Bryan Kaufman

\end{document}
